% BHU PYQ -- Manually transcribed & error-corrected
% Paper: MTB-301 : Algebra-I
% Exam: B.Sc./B.A. (Hons.) Semester III Examination, 2023-24

\documentclass[11pt,a4paper]{article}
\usepackage[a4paper,top=2.5cm,bottom=2.5cm,left=2.8cm,right=2.8cm,headheight=14pt]{geometry}
\usepackage{amsmath,amssymb}
\usepackage[T1]{fontenc}\usepackage[utf8]{inputenc}
\usepackage{lmodern,microtype,fancyhdr,enumitem,hyperref,lastpage,parskip}

\hypersetup{colorlinks=true,urlcolor=black,linkcolor=black,
  pdftitle={MTB-301 Algebra-I -- BHU BSc/BA Sem III 2023-24}}
\pagestyle{fancy}\fancyhf{}
\renewcommand{\headrulewidth}{0.4pt}\renewcommand{\footrulewidth}{0.4pt}
\fancyhead[L]{\small\textit{Banaras Hindu University}}
\fancyhead[R]{\small\textit{MTB-301: Algebra-I}}
\fancyfoot[C]{\small Page \thepage\ of \pageref{LastPage}}

\newlist{parts}{enumerate}{2}
\setlist[parts,1]{label=(\alph*),leftmargin=*,itemsep=5pt,topsep=3pt}
\setlist[parts,2]{label=(\roman*),leftmargin=*,itemsep=3pt,topsep=2pt}
\newcommand{\pts}[1]{\hfill\textit{[#1]}}

\begin{document}
\begin{center}
  {\large\bfseries BANARAS HINDU UNIVERSITY}\\[2pt]
  {\normalsize B.Sc./B.A.\ (Hons.) --- Semester III Examination, 2023-24}\\[6pt]
  \rule{0.8\linewidth}{0.5pt}\\[6pt]
  {\large\bfseries Paper No.\ MTB-301: Algebra-I}\\[4pt]
  {\normalsize Time: 3 Hours \hfill Full Marks: 70}\\[2pt]
  {\small Ref.\ No.:\ 067771}
\end{center}
\rule{\linewidth}{0.4pt}
\smallskip
\noindent\textit{Instructions: Answer \textbf{five} questions including Question~1 which is compulsory. Figures in right-hand margin indicate marks.}
\bigskip

\noindent\textbf{Question 1.} \textit{(Compulsory --- 2 marks each.)}
\begin{parts}
  \item $A$ is a real non-symmetric matrix of order 3. Find the rank of $A-A^T$.
  \item Give an example each of a Hermitian and a skew-Hermitian matrix of order 3.
  \item Find the value of $\alpha$ for which $x+y+z=1$, $x+2y-z=0$, $5x+7y+\alpha z=0$ has a unique solution.
  \item Find the number of elements of order 5 in $(\mathbb{Z}_{30},+)$.
  \item Define automorphism and give an example.
  \item Let $G=\{A\in M_{n\times n}:\det(A)\neq0\}$ and $H=\{B\in M_{n\times n}:\det(B)=1\}$. Prove $(H,*)$ is a normal subgroup of $(G,*)$.
  \item ``Every commutative group is cyclic.'' Justify.
  \item Find all right cosets of $(5\mathbb{Z},+)$ in $(\mathbb{Z},+)$.
  \item $G$ is cyclic of order 20 generated by $a$. Find the order of the subgroup generated by $a^5$.
  \item Let $G=\{x,\,1/x,\,-x,\,-1/x\,;\,x\in\mathbb{R}\setminus\{0\}\}$ with composition of functions. Find the identity and inverse of each element.
  \item Is there an isomorphism from $(\mathbb{Z}_4,+)$ to Klein's 4-group? Give reasons.
\end{parts}
\medskip\rule{\linewidth}{0.2pt}

\medskip\noindent\textbf{Question 2.}
\begin{parts}
  \item Prove that a complex matrix can be uniquely expressed as a sum of a Hermitian and a skew-Hermitian matrix.\pts{4}
  \item For an $n\times n$ complex matrix $A$, show $A^TA=O\Rightarrow A=O$.\pts{4}
  \item $A$ is real orthogonal and $I+A$ is non-singular. Prove $(I+A)^{-1}(I-A)$ is skew-symmetric.\pts{4}
\end{parts}
\medskip\rule{\linewidth}{0.2pt}

\medskip\noindent\textbf{Question 3.}
\begin{parts}
  \item Let $A=\begin{pmatrix}0&-1&0&0\\1&0&0&0\\0&0&0&-1\\0&0&1&0\end{pmatrix}$. Show $A^4=I$ and find $A^{51}$.\pts{6}
  \item Define rank of a matrix. Find the rank of $\begin{pmatrix}0&1&2&4\\1&3&4&8\\2&6&0&-4\\2&3&4&2\\-3&4&2&-3\end{pmatrix}$ by reducing to echelon form.\pts{6}
\end{parts}
\medskip\rule{\linewidth}{0.2pt}

\medskip\noindent\textbf{Question 4.}
\begin{parts}
  \item Define odd and even permutations. Is $S_3$ simple?\pts{4}
  \item Prove that arbitrary intersection of subgroups is a subgroup. Is this true for union? Justify.\pts{4}
  \item Show the centre $Z(G)$ of a group $G$ is a normal subgroup.\pts{4}
\end{parts}
\medskip\rule{\linewidth}{0.2pt}

\medskip\noindent\textbf{Question 5.}
\begin{parts}
  \item State and prove Cayley's theorem in group theory.\pts{7}
  \item Find the permutation group isomorphic to $G=(\{1,-1,i,-i\},\cdot)$.\pts{7}
\end{parts}
\medskip\rule{\linewidth}{0.2pt}

\medskip\noindent\textbf{Question 6.}
\begin{parts}
  \item State the fundamental theorem of homomorphism. Define $\varphi:(\mathbb{R},+)\to(\{z\in\mathbb{C}:|z|=1\},\cdot)$ by $\varphi(x)=e^{2\pi ix}$. Show $\varphi$ is a homomorphism, find $\ker\varphi$, and identify the quotient group.\pts{6}
  \item Prove every subgroup of a cyclic group is cyclic. Verify with an example.\pts{6}
\end{parts}
\medskip\rule{\linewidth}{0.2pt}

\medskip\noindent\textbf{Question 7.}
\begin{parts}
  \item Let $D$ be the set of all odd integers with $a*b=a+b-1$. Check whether $(D,*)$ is a group and if it is commutative.\pts{6}
  \item In a group $(G,\circ)$, show $(a\circ b)^{-1}=b^{-1}\circ a^{-1}$ for all $a,b\in G$.\pts{3}
  \item If $G$ is cyclic and $G=\langle a\rangle$, show $o(G)=o(a)$.\pts{3}
\end{parts}

\vfill\rule{\linewidth}{0.4pt}
\begin{center}\small
\href{https://bhu-exam.vercel.app/}{Click here to visit our website}\\[4pt]
\href{https://me-aryan.vercel.app/}{Click here to visit our contributor}\\[4pt]
\href{https://quashan-me.vercel.app/}{Click here to visit our contributor}
\end{center}
\end{document}
